\subsection{Work Remaining (Layer IV: Search- and Documentation-Level Poisoning of Coding Agents)}

\noindent \textbf{Position in the thesis.}
This remaining work is the fourth layer of the dissertation.
Layers I and II focus on smart-contract runtime defense (invariants and control-flow integrity),
Layer III focuses on auditing harmful artifacts in LLM-generated code,
and Layer IV extends the off-chain line to an upstream channel:
search- and documentation-level poisoning of coding agents.

\paragraph{Motivation.}
Modern coding agents increasingly rely on external documentation, tutorials, issue threads,
and search results while solving software-engineering tasks.
This improves utility but introduces a new attack surface:
attackers may influence agent behavior by controlling external content that is retrieved and trusted,
without modifying model weights, system prompts, or repository code.
The core practical question is end-to-end:
when poison is visible in retrieval channels, how often is it retrieved,
accepted by the agent, and translated into harmful code actions?

\paragraph{Problem statement.}
This work studies \emph{external-content poisoning against coding agents}.
The attacker controls a subset of externally visible content (e.g., documentation pages,
API references, tutorial fragments, issue discussions, or search snippets),
while the victim is a coding agent that retrieves such content during task solving.
The objective is to induce harmful outcomes such as insecure patches,
plausible but incorrect patches, or unsafe recommendations.
The central challenge is to measure and explain compromise across the full pipeline:
retrieval, trust, and action.


\paragraph{Research questions.}
\begin{itemize}
    \item \textbf{RQ1 (Baseline vulnerability).}
    To what extent are open-source coding agents vulnerable to poisoned search results
    and documentation pages during realistic software-engineering tasks?
    \item \textbf{RQ2 (Budget/rank/stealth).}
    How do success rates vary with attacker budget, retrieval rank,
    corroborating poisoned pages, and detectability constraints?
    \item \textbf{RQ3 (Defense tradeoff).}
    Which practical defenses reduce compromise while preserving benign coding performance?
\end{itemize}

\paragraph{Threat model.}
We assume a black-box attacker who cannot alter model weights,
agent system prompts, benchmark labels, or repository state,
but can control some external search-visible and documentation-level content.
The victim is a coding agent that can browse/read external information while solving tasks.
Primary outcomes of interest are:
(1) incorrect but plausible patches,
(2) security-degrading patches,
and (3) unsafe recommendations or citations in final outputs.

\paragraph{Experimental design.}
The empirical study will evaluate at least two open-source coding-agent stacks
(e.g., OpenHands-style and SWE-agent-style systems)
using task suites that require external lookup.
To ensure reproducibility, the main benchmark will use a fixed local corpus
combining benign and attacker-controlled documents,
with a smaller live-web case study as supplementary evidence.
Attack variants include:
rank poisoning, content poisoning, consensus poisoning,
and snippet/title poisoning.

\paragraph{Defenses and expected deliverables.}
The planned defense study will compare lightweight and deployable mechanisms:
source-provenance filtering, instruction/content separation,
cross-source agreement gating, and post-generation security checking.
Expected deliverables are:
\begin{itemize}
    \item a reproducible benchmark for search/documentation poisoning of coding agents,
    \item a stage-aware empirical analysis of compromise pathways,
    \item a security--utility characterization of practical defenses, and
    \item deployment-oriented guidance for reducing upstream AI-assisted development risk in Web3 pipelines.
\end{itemize}

\paragraph{Scope control.}
To keep this layer focused and publishable,
the work will target coding agents and documentation/search channels only;
it will not attempt to cover all browser-agent environments or all general prompt-injection settings.
This scope aligns Layer IV with the thesis objective of practical,
engineering-relevant security hardening.
